\section{Hyperparameters}
\label{app:hyper}

Table~\ref{tab:hyper} lists the full hyperparameter set used to
produce the deployed model. All values are unchanged across the
five training seeds.

\begin{table}[h]
\centering
\caption{Training and deployment hyperparameters for the deployed
\fastgrnn{} model.}
\label{tab:hyper}
\resizebox{\columnwidth}{!}{%
\begin{tabular}{lr}
\toprule
Parameter & Value \\
\midrule
Hidden size $H$                  & 16 \\
Input dimensionality $d$         & 3 (tri-axial accel.) \\
Window length $T$                & 128 samples (2.56~s) \\
Sample rate                      & 50~Hz \\
Output classes                   & 6 \\
\midrule
Recurrent rank $r_u$             & 8 \\
Input rank $r_w$                 & 2 \\
Target sparsity $s$              & 0.5 (181 cell $+$ 102 head $=$ 283 nonzero / 566~B) \\
Sparsity ramp epochs             & 50 (cubic schedule) \\
Frozen-mask fine-tune epochs     & 50 \\
\midrule
Optimizer                        & Adam \\
Learning rate $\eta$             & $10^{-3}$ \\
Batch size                       & 64 \\
Total epochs                     & 100 \\
Random seeds                     & $\{0, 1, 2, 3, 4\}$ \\
\midrule
Activation function              & sigmoid, tanh \\
LUT size                         & 256 entries each \\
LUT input domain                 & $[-8, +8]$ \\
Weight quantization              & per-tensor Q15 \\
Activation quantization          & per-tensor Q15 + calibration \\
Calibration mini-batches         & 5 \\
Calibration headroom             & 10\% \\
\bottomrule
\end{tabular}}
\end{table}

\section{Q15 Quantization Implementation}
\label{app:q15}

The per-tensor Q15 quantization step used in
Section~\ref{sec:method:q15} is a one-line operation per tensor. In
Python:

\begin{verbatim}
def to_q15(W: np.ndarray, scale: float) -> np.ndarray:
    return (W / scale).round() \
                      .clip(-32768, 32767) \
                      .astype(np.int16)
\end{verbatim}

The per-tensor scale is chosen as
\[
s_\ell = \max_{ij}\bigl|W^{(\ell)}_{ij}\bigr| \,\big/\, 32767,
\]
which is the largest scale that keeps every quantized entry inside
the signed 16-bit range. The same expression is used for both the
recurrent weight factors $\mathbf{W}_1, \mathbf{W}_2,
\mathbf{U}_1, \mathbf{U}_2$ and the classifier head. At runtime, the
dequantized value is computed as

\begin{verbatim}
float w = (float)W_q15[i] * scale;
\end{verbatim}

The runtime cost per dequantize is one int-to-float conversion and
one multiply; both are cheap on the AVR and, with the LUT
intercepting the activations, manageable on the multiplier-less
\msp{}.

\section{LUT Generation Algorithm}
\label{app:lut}

The sigmoid and tanh look-up tables of Section~\ref{sec:method:lut}
are pre-computed offline in Python and emitted as a C header:

\begin{verbatim}
LUT_SIZE     = 256
INPUT_MIN    = -8.0
INPUT_MAX    =  8.0
BUCKET_WIDTH = (INPUT_MAX - INPUT_MIN) / LUT_SIZE

sigmoid_lut = [
    1.0 / (1.0 + math.exp(-(INPUT_MIN
            + (i + 0.5) * BUCKET_WIDTH)))
    for i in range(LUT_SIZE)
]
tanh_lut = [
    math.tanh(INPUT_MIN + (i + 0.5) * BUCKET_WIDTH)
    for i in range(LUT_SIZE)
]
\end{verbatim}

Each entry stores the value of the activation at the
\emph{center} of its bucket (the $(i + 0.5)$ offset), which is
the maximum-likelihood estimate for a uniformly distributed
sub-bucket input and avoids a half-bucket bias. The runtime lookup is

\begin{verbatim}
static inline float lut_eval(const float *lut, float x) {
    if (x <= LUT_INPUT_MIN) return lut[0];
    if (x >= LUT_INPUT_MAX) return lut[LUT_SIZE - 1];
    int idx = (int)((x - LUT_INPUT_MIN) * LUT_INPUT_SCALE);
    return lut[idx];
}
\end{verbatim}

with \texttt{LUT\_INPUT\_SCALE = 1 / BUCKET\_WIDTH}. The two tables
together occupy 2~KB of Flash; together they eliminate every
\texttt{expf} and \texttt{tanhf} call from the runtime.

\section{MSP430 I\textsuperscript{2}C Bring-Up}
\label{app:hw}

Live MPU-6050 read-out on the \msp{} uses the USCI\_B0 peripheral
in I\textsuperscript{2}C master mode with the standard pin
assignment (P1.6~$=$~SCL, P1.7~$=$~SDA). Two practical issues are
worth recording for future readers:

\begin{itemize}
\item \textbf{J5 jumper removal.} The MSP-EXP430G2 LaunchPad
multiplexes the on-board green LED with P1.6. The J5 jumper must
be physically removed before any I\textsuperscript{2}C
communication; otherwise SCL is loaded by the LED and the bus
sits below the I\textsuperscript{2}C high-level threshold.

\item \textbf{Bus-busy recovery.} The USCI\_B0 peripheral can latch
into a \texttt{UCSCLLOW + UCBBUSY} state if power is removed
mid-transfer; the conventional 9-clock recovery sequence
(toggle SCL nine times with SDA released, then STOP) restores the
bus. A GPIO-level health check at boot
(\texttt{P1IN \& (BIT6|BIT7)} should be \texttt{0xC0}) detects the
condition before the I\textsuperscript{2}C state machine is engaged.
\end{itemize}

These two items consumed approximately one day of bring-up time
and are not documented in the standard MSP430 application
notes~\cite{ti_msp430}. They are included here in case they save
the next reader the same day.
